\section{Literature Review}

\subsection{Query Expansion in RAG System}

Query expansion (QE) improves retrieval performance by simulating PRF. With the rise of LLMs, recent studies have explored LLMs as zero-shot QE engines. HyDE~\cite{gao2023precise} and Query2Doc~\cite{wang2023query2doc} demonstrate that prompting LLMs can yield precise expansions with PRF. Other works enhance semantic diversity or reasoning depth: MILL~\cite{jia2023mill} employ mutual verification, and GRF~\cite{mackie2023generative} apply LLM-based relevance feedback. ThinkQE~\cite{lei2025thinkqe} models expansion as a reasoning process, while LLMlingua~\cite{jiang2023llmlingua} focuses on compressing expansion prompts. GAR~\cite{xia2024knowledge} further introduces a knowledge graph-aware QE framework.

However, existing work overlooks personalization and fails to align QE with user-specific semantic spaces. In contrast, our PBR generates style-aware and structure-aware expansions, bridging this gap for personalized retrieval.

\subsection{Personalized RAG System}

Personalized RAG system has recently emerged as a crucial direction for adapting LLMs to individual users. A recent survey by Li et al.~\cite{li2025survey,li2023agent4ranking} provides a comprehensive overview, highlighting the importance of personalization in RAG. Approaches such as UniMS-RAG~\cite{wang2024unims} integrate multi-source signals for dialogue personalization, while PersonaRAG~\cite{zerhoudi2024personarag} enhances personalization via role-based agent frameworks. Salemi et al.~\cite{salemi2024lamp} further investigate LLM personalization by modeling language preferences and user goals. Complementary work like Wu et al.~\cite{wu2024understanding} explores the role of explicit user profiles, and Cohn et al.~\cite{cohn2025personalizing} leverage interaction logs for educational agent adaptation. Domain-specific applications, such as personalized care systems~\cite{yang2025improving} and memory-based assistant evaluation~\cite{wu2025longmemeval,xu2025towards,xu2025align}, highlight the importance of long-term memory and evolving user states in effective RAG.

Despite recent advances, little attention has been given to improving retrieval accuracy at the pre-retrieval stage. PBR enables fine-grained personalization by jointly modeling the user’s style and corpus structure before retrieval.


% \subsection{Query Expansion}

% Query expansion (QE) has long been recognized as a critical technique to enhance retrieval effectiveness by reformulating user queries into richer semantic representations. With the advent of large language models (LLMs), a growing body of work has explored how LLMs can serve as powerful engines for zero-shot or few-shot QE.

% Recent studies have demonstrated the capacity of LLMs to generate informative expansions without the need for explicit supervision. Gao et al.~\cite{gao2023precise} propose a precise zero-shot dense retrieval method that avoids reliance on relevance labels, instead leveraging LLMs to synthesize queries that align closely with target document semantics. Jagerman et al.~\cite{jagerman2023query} further explore prompting strategies that guide LLMs toward producing high-quality expansions, showing that simple prompt templates can outperform traditional QE pipelines.

% Incorporating mutual verification, Jia et al.~\cite{jia2023mill} introduce MILL, which employs dual LLM agents to validate the consistency and informativeness of generated expansions. Wang et al.~\cite{wang2023query2doc} present Query2Doc, a generation-based approach that reimagines queries as pseudo-documents to bridge semantic granularity gaps in retrieval. Meanwhile, Mackie et al.~\cite{mackie2023generative} propose using generative relevance feedback from LLMs to iteratively refine query intent based on initial retrieval results.

% More recently, Lei et al.~\cite{lei2025thinkqe} advocate for a thinking-style QE process in ThinkQE, wherein LLMs simulate a chain-of-thought reasoning trajectory to incrementally evolve query representations. In a related direction, Jiang et al.~\cite{jiang2023llmlingua} investigate prompt compression techniques, such as LLMlingua, to reduce query redundancy while preserving critical semantics—an idea indirectly beneficial to compact and efficient QE.

% Overall, these methods highlight the growing shift from static expansion rules to dynamic, reasoning-driven query augmentation, laying the groundwork for more personalized and context-aware retrieval systems.